\documentclass[11pt]{article}
\usepackage{fontspec}
\usepackage[french,english]{babel}
\usepackage{amsmath}
\usepackage{amssymb}
\usepackage{amsthm}
\usepackage{mathrsfs}
\usepackage{tikz-cd}
\usepackage{enumitem}

\newcommand{\printedpage}[1]{\marginpar{\footnotesize p.~#1}}
\newcommand{\msp}[1]{\marginpar{\footnotesize [#1]}}
\newcommand{\Ch}{\underline{\mathrm{Ch}}}
\newcommand{\Ty}{\mathrm{Ty}}
\newcommand{\uTy}{\underline{\mathrm{Ty}}}
\newcommand{\Fl}{\mathrm{Fl}}
\newcommand{\Hom}{\mathrm{Hom}}
\newcommand{\ch}{\mathrm{ch}}
\newcommand{\hot}{\mathrm{hot}}
\newcommand{\id}{\mathrm{id}}

\begin{document}

\printedpage{44}
% head: A. Grothendieck : Dérivateurs, Ch. VII, Édition des Universités de Montpellier II et Paris VII
% foot: Version 22 décembre 2001

which is in accord with the intuition that the operation $\gamma \longmapsto \gamma^o$ plays the role of passing to the inverse (``up to homotopy''), and we have the formula
\[
(\gamma'\gamma)^{-1} = \gamma^{-1}\gamma'^{-1}
\]
for two invertible arrows $\gamma$, $\gamma'$ composable (in the order $\gamma'$, $\gamma$).

As for the universal property of
\[
X \longrightarrow \Pi(X)
\]
(which will allow us to write
\[
\Pi(X) \simeq X\Sigma^{-1}, \qquad \text{with } \Sigma = \Fl(X) \ \ ),
\]
we may as well establish it right away in the general setting where we give ourselves a subset
\[
\Sigma \subseteq \Fl(X),
\]
and we define, for $\tau \in \Ty$,
\[
\Ch_{\Sigma,\tau}(X) \subseteq \Ch_{\tau}(X)
\]
as the full subcategory of $\Ch_{\tau}(X)$ consisting of the paths in $X$ for which \textit{all the retrograde arrows}

\msp{page 66}

are in $X$ \texttt{[rather in $\Sigma$]}. We shall note that for a path $c$ and its opposite $c^o$ both to be $\Sigma$-\textit{paths,} it is necessary and sufficient that \textit{all} the arrows of the path $c$ be in $\Sigma$ (we shall then say that $c$ is a $\Sigma$-\textit{strict path}), since the retrograde arrows of $c^o$ are the direct arrows of $c$.

For the full subcategories $\Ch_{\tau,\Sigma}(X)$ to be stable under the inverse image functors
\[
f^* : \Ch_{\tau'}(X) \longrightarrow \Ch_{\tau}(X)
\]
associated with the arrows
\[
f : \tau \longrightarrow \tau'
\]
in $\uTy$, or simply for there to be stability for the two functors $f^*$ associated with an $f : \Delta_- \longrightarrow 1$ and $\Delta_- \longrightarrow \Delta_-^2$, it is necessary and sufficient that $\Sigma$ be stable under composition and contain all the identity arrows % [sic] qu'y ait -> qu'il y ait, missing word
\[
\left\{
\begin{aligned}
& x \in \mathrm{Ob}\,X \ \Longrightarrow \ \id_x \in \Sigma \\
& u, v \in \Sigma, \ \ u,\, v \text{ composable} \Longrightarrow uv \in \Sigma.
\end{aligned}
\right.
\]
Indeed, the arrows in a path $f^*(c') \in \Ch_{\tau}(X)$ are either identity arrows (which correspond to the transition arrows of $I_\tau$ degenerate for $f$), or arrows among those of $\tau$, or composites of a sequence of consecutive composable arrows of $\tau'$ (which correspond to the transition arrows in $I_{\tau'}$ included in the part of a transition arrow of $I_\tau$), with preservation of signatures.

\printedpage{45}
% head: A. Grothendieck : Dérivateurs, Ch. VII, Édition des Universités de Montpellier II et Paris VII
% foot: Version 22 décembre 2001

\msp{page 67}

We assume these conditions satisfied, so that the $\Ch_{\tau,\Sigma}(X)$ form the fibers of a full fibered subcategory of $\Ch(X)$ over $\uTy$, which we shall denote by
\[
\Ch_{\Sigma}(X) \hookrightarrow \Ch(X).
\]
It is obvious that this subcategory is stable under the operation of composition of paths (but not, in general, as we have said, under passage from $c$ to $c^o$ (for this $\Sigma$ would have to contain \textit{all} the arrows of $X$, since for every $u : a \longrightarrow b$, which is thus in $\Ch_{\Delta_+,\Sigma}(X)$, $u^o$ given by $b \xleftarrow{\ u\ } a$ in $\Ch_{\Delta_-}(X)$ would have to be a $\Sigma$-path, whence $u \in \Sigma$)).

We are going to define a canonical category
\[
\Pi_\Sigma(X),
\]
having the same objects as $X$, and a canonical functor
\[
X \longrightarrow \Pi_\Sigma(X)
\]
which is the identity on objects, and which transforms the $u \in \Sigma$ into invertible arrows. For this we copy the description of the fundamental groupoid (but $\Pi_\Sigma(X)$ will not in general be a groupoid), by putting
\[
\Hom_\Sigma(x,y) = \pi_0(\Ch_\Sigma(x,y)),
\]

\msp{page 68}

and defining the composition of ``$\Sigma$-arrows'' in terms of the composition of paths. The argument which proved, in the case of $\Pi(X)$, that \textit{every} arrow of $X$, viewed as a path $c$ of type $\Delta_+$, was invertible in $\Pi(X)$, having as inverse $c^o$ of type $\Delta_-$, is still valid, provided that $c^o$ is indeed still a $\Sigma$-path, which requires that $u \in \Sigma$.

We are going to interpret it in terms of a canonical functor
\[
\kappa : X \longrightarrow \Pi_\Sigma(X) \qquad \text{(\footnotemark[19])}
\tag*{$(*)$}
\]
\footnotetext[19]{On arrows, $\kappa$ is defined, for $u : a \longrightarrow b$, by $\kappa(u) =$ homotopy class of $\ch(u)$ in $\Ch_\Sigma(a,b)$, where $\ch(u)$ is the path of type $\Delta_+$ from $a$ to $b$ defined by $u$.}
which is again described as in the case of the fundamental groupoid, by associating to every arrow $u$ of $X$ the homotopy class of the path $c(u)$ it defines. We thus see that this canonical functor transforms the $u \in \Sigma$ into invertible arrows.

When the arrows of $\Sigma$ are already invertible in $X$, I say that the functor $(*)$ is an isomorphism, i.e. that for every $a, b \in \mathrm{Ob}\,X$,
\[
\alpha : \Hom_X(a,b) \longrightarrow \pi_0\Ch_\Sigma(a,b), \qquad u \longmapsto \underbrace{\ch_\Sigma(u)}_{= \ch(u)}
\]
is bijective. To see this, we are going to define
\[
\beta : \pi_0(\Ch_\Sigma(a,b)) \longrightarrow \Hom_X(a,b),
\]
which will be inverse to $\beta$. We define % [sic] inverse de $\beta$ -> inverse de $\alpha$, $\beta$ is meant to be the inverse of $\alpha$

\printedpage{46}
% head: A. Grothendieck : Dérivateurs, Ch. VII, Édition des Universités de Montpellier II et Paris VII
% foot: Version 22 décembre 2001

\msp{page 69}

$\beta(c)$, for a $\Sigma$-arrow of type $\tau$, the length $n$, by inverting all the retrograde arrows (which is possible, since they are in $\Sigma$), so as to obtain a path $c^\natural$ of type $\Delta_+^n$, and by taking the composite (in $X$) of the arrows of this path $c^\natural$. It is then obvious that $\beta\alpha = \id$, we must prove that we also have $\alpha\beta = \id$, i.e. that in $\Pi_\Sigma(X)$, the arrows $u = \int c^\natural$ and $c$ are equal, i.e. $u$ and $c$ are homotopic in $\Ch_\Sigma(a,b)$. Let us consider the path $c^\natural$, of type $\Delta_+^n$: % [sic] le longueur -> la longueur, gender; possibles -> possible, number
\[
a = a_0 \xrightarrow{u_1} a_1 \xrightarrow{u_2} a_2 \longrightarrow \ \cdots \ \longrightarrow a_{n-1} \xrightarrow{u_n} a_n = b,
\]
and let $u_i$ ($i \in [1,n]$) be an arrow of this path, corresponding to a transition arrow $\underline{u}_i$ that is \textit{retrograde} in $\tau$, hence by construction
\[
u_i = v_i^{-1} \qquad \text{in } X\text{, with } v_i \in \Sigma
\]
the arrow with index $i$ in the initial path $c$ of type $\tau$. If $\underline{u}_i$, on the other hand, is direct, we will have
\[
u_i = v_i.
\]
(NB We denote by $v_i$ the $i$-arrow in the path $c$.) Now we have
\[
c = c_n * c_{n-1} * \cdots * c_1 \qquad \text{in } \Pi_\Sigma(X)
\tag*{$(*)$}
\]
and the factor $c_i$ is either the homotopy class of $\ch(v_i) = \ch(u_i)$ (the case where the transition arrow $\underline{u}_i$ of $\tau$ is direct), or that of the path \textit{opposite} to the path $\ch(v_i) : a_i \xrightarrow{v_i} a_{i-1}$ of type $\Delta_+$, hence the path $\ch(v_i)^o : a_{i-1} \xleftarrow{v_i} a_i$ of type $\Delta_-$. I say that
\[
\boxed{\ \ch(v_i)^o = \ch(u_i) \qquad \text{\textit{in} } \Pi(X),\ }
\tag*{(1)}
\]
which proves that in $\Pi_\Sigma(X)$, $c$ is equal to the composite of the $\ch(u_i)$ ($i \in [1,n]$), hence to $c^\natural$. We must prove moreover that
\[
c^\natural = \ch(u) \qquad (u = u_nu_{n-1}\cdots u_1) \qquad \text{in } \Pi(X).
\tag*{(2)}
\]

\msp{page 70}

So there are two verifications to make (where I again assume $\Sigma$ arbitrary):

\begin{enumerate}[label={}]
\item[$1^o$)] Let $u \in \Sigma$, $u : a \longrightarrow b$, suppose $\Sigma$ invertible in $X$. Then
\[
\ch(u)^o = \ch(u^{-1}) \qquad \text{in } \Pi_\Sigma(X),
\]
i.e. the two paths from $b$ to
\[
\begin{array}{cccl}
b & \xleftarrow{\ u\ } & a & \text{of type } \Delta_- \\
b & \xrightarrow{u^{-1}} & a & \text{of type } \Delta_+
\end{array}
\]
\printedpage{47}
% head: A. Grothendieck : Dérivateurs, Ch. VII, Édition des Universités de Montpellier II et Paris VII
% foot: Version 22 décembre 2001
are homotopic in $\Ch_\Sigma(b,a)$. Now we already know, since $\kappa : X \longrightarrow \Pi_\Sigma(X)$ is a functor, that
\[
\hot\,\ch(u^{-1}) = (\hot\,\ch(u))^{-1} \qquad \text{in } \Pi_\Sigma(X)
\]
where $\hot(\bullet)$ denotes the homotopy class of a chain (here a class in $\Ch_\Sigma(b,a)$ and in $\Ch_\Sigma(a,b)$ respectively). So the statement to prove is that $\hot(\ch(u)^o)$ is equal to the \textit{inverse} of $\hot(\ch(u))$ in $\Pi_\Sigma(X)$. This is thus what we proved for \textit{every} $u \in \Fl(X)$, in the case of $\Pi(X)$. The argument must work, so long as it makes sense, i.e. that $\ch(u)^o$ really is a $\Sigma$-path, i.e. $u \in \Sigma$.

\item[$2^o$)] Formula (2) means that for every path $c$ of type $\Delta_+^n$ (thus tautologically a \msp{page 71} $\Sigma$-path), going from $a$ to $b$, its class in $\Hom_\Sigma(a,b)$ is none other than that of $\ch(u)$, where $u$ is the composite $u_nu_{n-1}\cdots u_1$. But we have
\[
c = c_1 * c_2 * \cdots * c_n,
\]
whence in $\Pi_\Sigma(X)$
\[
\hot(c) = \hot(c_n)\hot(c_{n-1})\cdots\hot(c_1),
\]
now $\hot(c_i) = \hot\,\ch(u_i) = \kappa(u_i)$, hence
\[
\begin{aligned}
\hot(c) \ &= \ \kappa(u_n)\kappa(u_{n-1})\cdots\kappa(u_1) \\
&= \ \kappa(\underbrace{u_nu_{n-1}\cdots u_1}_{u}),
\end{aligned}
\]
since $\kappa$ is a functor, hence
\[
\hot(c) = \hot\,\ch(u),
\]

\hfill q.e.d.
\end{enumerate}

Let us now note that if $(X', \Sigma')$ is a second pair consisting of a category $X'$ and an \textit{``admissible''} set of arrows of $X'$, and if
\[
f : X \longrightarrow X'
\]
is a functor such that
\[
f(\Sigma) \subseteq \Sigma',
\]
then there exists a \textit{unique} functor
\[
\bar{f} = \Pi_\Sigma(f) : \Pi_\Sigma(X) \longrightarrow \Pi_{\Sigma'}(X')
\]
which ``extends'' $f$, i.e. such that the square
\printedpage{48}
% head: A. Grothendieck : Dérivateurs, Ch. VII, Édition des Universités de Montpellier II et Paris VII
% foot: Version 22 décembre 2001
\msp{page 72}
\[
\begin{tikzcd}
X \arrow[r, "\kappa_X"] \arrow[d, "{\scriptstyle f}"] & \Pi_\Sigma(X) \arrow[d, "{\scriptstyle \bar{f} = \Pi_\Sigma(f)}"] \\
X' \arrow[r, "\kappa_{X'}"] & \Pi_{\Sigma'}(X')
\end{tikzcd}
\]
is commutative. Uniqueness follows from the fact that on objects $\bar{f}$ must be given by $f$, and that on arrows the commutativity of the square implies that $\bar{f}$ is known on arrows of the form $\kappa_X(u)$, and thus also on the inverses of such arrows, when these inverses are invertible, hence in particular on the $\kappa_X(u)^{-1}$ for $u \in \Sigma$. Now every arrow of $\Pi_\Sigma(X)$, coming from a path, is a composite of a sequence of arrows of this form. For the existence of $\bar{f}$, we note that $f$ induces functors % [sic] les flèche -> les flèches, plural
\[
\Ch_\tau(X; a, b) \longrightarrow \Ch_\tau(X'; f(a), f(b))
\]
which commute with the operations $\varphi^*$ on paths associated with the functors
\[
\varphi : \tau \longrightarrow \tau',
\]
whence a homomorphism of fibered categories over $\uTy$
\[
f_{a,b} : \Ch(X; a, b) \longrightarrow \Ch(X'; f(a), f(b)).
\]

\msp{page 73}

Whence a map on the $\pi_0$, i.e.
\[
\bar{f}_{a,b} : \Hom_\Sigma(a,b) \longrightarrow \Hom_\Sigma(f(a), f(b)),
\]
given by
\[
\bar{f}_{a,b}(\hot(c)) = \hot(f_{a,b}(c)).
\]
We obviously have, for two paths $\gamma : a \longrightarrow b$ and $\gamma' : b \longrightarrow c$,
\[
f_{a,c}(\gamma\gamma') = f_{b,c}(\gamma')f_{a,b}(\gamma),
\]
which proves that $\bar{f}$ is compatible with composition. It is trivial that it is compatible with identities (for $f_{a,a}$ transforms the empty path at $a$ into the empty path at $f(a)$), so we have a functor $\bar{f}$, and it really does make the square commutative.

We can now prove the universal property of $\Pi_\Sigma(X)$ and of
\[
X \xrightarrow{\ \kappa\ } \Pi_\Sigma(X),
\]
namely that every functor $f : X \longrightarrow X'$ such that $u \in \Sigma \Longrightarrow f(u)$ invertible, factors \textit{uniquely} through $\Pi_\Sigma(X)$. Uniqueness is seen as before (\footnote[20]{This is unnecessary, cf. below.}). For existence, we note that $f$ maps $\Sigma$ into the set $\Sigma'$

\end{document}
